\documentclass[../main.tex]{subfiles}

\begin{document}

% Hvert bilag er en \section. I appendices-miljøet i main.tex bliver de til
% "Bilag A", "Bilag B" osv., også i indholdsfortegnelsen.
% Henvis til et bilag i teksten med \cref{app:eksempel} -> "bilag A".
%
% PDF-filer lægges i denne mappe (Bilag/) og indsættes med \includepdf.
% \subfix{...} gør, at stien også virker, når bilaget kompileres for sig selv.
%
% Et PDF-bilag med flere sider indsættes sådan:
%   \includepdf[pages=1, scale=0.8, pagecommand={\section{Navn}\label{app:navn}}]{\subfix{fil.pdf}}
%   \includepdf[pages=2-, scale=0.8, pagecommand={}]{\subfix{fil.pdf}}
% Første linje sætter bilagets overskrift oven på PDF'ens første side, anden linje
% tager resten af siderne (slet den, hvis PDF'en kun har én side).
% scale=0.8 giver plads til headeren og sidetallet.

% --- Eksempel ---
\section{Eksempel}
\label{app:eksempel}
\todo[inline]{Læg bilag som PDF i Bilag/ og indsæt dem med \textbackslash includepdf (se kommentaren øverst)}

\end{document}
